Large language models generate plausible but factually incorrect outputs, motivating uncertainty-based hallucination detection. Token entropy and semantic consistency are two promising signals, yet their relative effectiveness and potential complementarity remain unstudied. We present the first systematic comparison of these signals on TriviaQA using Llama-2-7B-chat. Our experiments reveal that semantic consistency---measured via pairwise embedding similarity across multiple samples---achieves AUROC 0.81 with large effect size (Cohen's d=1.19), dramatically outperforming token entropy (AUROC 0.65, d=0.47). Surprisingly, linear fusion of these signals yields no improvement over consistency alone at proof-of-concept scale (N=20), with optimal weights collapsing to pure consistency. The moderate negative correlation between signals (r=-0.54) suggests partial redundancy rather than orthogonality. Our findings indicate practitioners should prioritize consistency-based detection for factual QA; the complementarity hypothesis requires larger-scale validation before practical deployment.
